%% Mentoring Plan
%% NSF 26-506 (PESOSE), Track 3 -- RedTwin
%%
%% Compile:  pdflatex mentoring_plan.tex   (or `make mentoring` from the repo root)
%%
%% COMPLIANCE NOTES (PAPPG 24-1, II.D.2.i(i)):
%%   - ONE PAGE MAXIMUM. Check the page count after every edit.
%%   - The plan must cover all postdoctoral scholars AND graduate students
%%     supported by the project, "regardless of whether they reside at the
%%     submitting organization, any subrecipient organization, or at any
%%     organization participating in a simultaneously submitted collaborative
%%     proposal." ASU is a subrecipient, so ASU must be covered; the previous
%%     draft described UCSB only.
%%   - Only ONE plan is submitted for the entire project.
%%   - Evaluated under the Broader Impacts review criterion.
%%   - PAPPG names these example activities: career counseling; training in
%%     preparation of proposals, publications and presentations; guidance on
%%     improving teaching and mentoring skills; guidance on collaborating with
%%     researchers from diverse backgrounds and disciplinary areas; and
%%     training in responsible professional practices. All five are addressed
%%     below; the previous draft omitted proposal preparation and responsible
%%     professional practices.
%%   - II.C.2 font/spacing/margin rules apply to supplementary documents.

\documentclass[11pt,letterpaper]{article}
\usepackage[letterpaper,margin=1in,heightrounded,
            headheight=0pt,headsep=0pt,footskip=0pt]{geometry}
\usepackage[T1]{fontenc}
\usepackage[utf8]{inputenc}
\usepackage{times}
\usepackage{microtype}
\usepackage{enumitem}
\usepackage{titlesec}
\usepackage{xcolor}

%% Conspicuous placeholder for facts only the PIs can confirm.
\newcommand{\todoverify}[1]{%
  \textbf{\textcolor{red}{[PI TO SUPPLY: #1]}}}

\setlength{\parindent}{0pt}
\setlength{\parskip}{0.45em}
\setlist[enumerate]{topsep=0.2em,itemsep=0.15em,parsep=0pt,leftmargin=1.6em}
\pagestyle{empty}

\begin{document}

\begin{center}
  {\large\bfseries Mentoring Plan}
\end{center}

\vspace{0.2em}

This plan covers every graduate student and postdoctoral researcher supported
by this project, at UC Santa Barbara and at Arizona State University, the
subrecipient. Mentoring has three components: individual mentoring by the PIs,
participation in departmental activities, and participation in
institution-level programs. The components apply to both groups, with the
differences noted below.

\textbf{Individual mentoring by the PIs.}
Every researcher funded by this project will be advised or co-advised by a PI
and will join that PI's laboratory. The four PIs have extensive experience in
graduate and postdoctoral mentoring and in placing their mentees in
high-quality positions in academia and industry. Each PI will provide:

\begin{enumerate}[label=\arabic*)]
  \item \emph{Career counseling.} Regular one-on-one meetings on career goals,
  progress, and strategy, and facilitated introductions to researchers,
  industry professionals, and alumni. The PIs will advise directly on the
  academic job search, on what to expect at an interview, and on the first
  years of an academic career. Postdoctoral researchers will additionally
  receive guidance on assembling a research statement and a teaching
  portfolio.
  \item \emph{Training in the preparation of proposals, publications, and
  presentations.} Mentorship in writing and submitting papers and rebuttals;
  training in presentation skills, including slide design and conveying
  results to non-specialist audiences; and, for senior students and
  postdoctoral researchers, participation in drafting sections of real
  proposals and in responding to reviews.
  \item \emph{Teaching and mentoring skills.} Opportunities to teach class
  sessions and to mentor junior researchers, including the undergraduates
  placed with the project through the Early Research Scholars Program at UCSB.
  \item \emph{Collaboration across disciplines and institutions.} The project
  is jointly executed by two universities and engages maintainers and
  operators outside academia, so researchers will collaborate with people from
  backgrounds and disciplines other than their own as a matter of course, with
  the PIs providing guidance on doing so effectively.
  \item \emph{Responsible professional practices.} Because this project
  discovers exploitable defects in widely deployed software, mentees will be
  trained in the coordinated vulnerability disclosure process the project
  follows, in the authorization and rules-of-engagement requirements that
  govern any assessment, and in the dual-use considerations that attach to
  offensive security research, alongside the standard responsible-conduct-of-research
  curriculum at their institution.
\end{enumerate}

Each mentee will maintain an Individual Development Plan, reviewed with their
advising PI at least annually and used to set the goals above against a
timeline.

\textbf{Departmental activities.}
Researchers will be encouraged to attend departmental mentoring events at both
institutions, including seminars on writing and presentation, effective
teaching techniques, and career development.

\textbf{Institution-level programs.}
At UCSB, the Center for Science and Engineering Partnerships (CSEP) runs an
intensive Professional Development Series covering basic, communication, and
career skills, in workshops of one to two hours offered roughly bi-weekly and
led by faculty and practitioners from STEM departments. 
At ASU, the mentoring of graduate students will take place through the Center
for Cybersecurity and Trusted Foundations.
The PIs will track participation and discuss it at the annual project meeting.

\end{document}
